\section*{Chapter \ref{ch:b2:ramsey-model} --- The Ramsey Model}

\begin{solution}{exo:b2:ramsey-model:1}
% figdata: saving_summary
Korea's rate rose from 1.0 per cent in 1960 to a peak of 41.7 in 1988, and was 34.0 in 2024. Japan's was already 44.6 in 1970, its highest, and fell to 26.8: Japan
was saving most when Korea was saving least.
\end{solution}

\begin{solution}{exo:b2:ramsey-model:2}
% code: exercises/solutions.py steady
$k^* = ((1/3)/0.08)^{1.5} = 8.505$; $y^* = 8.505^{1/3} = 2.041$; $c^* = 2.041 - 0.06 \times 8.505 = 1.531$; saving rate $0.06 \times 8.505/2.041 = 0.25$.
\end{solution}

\begin{solution}{exo:b2:ramsey-model:3}
% figdata: ky_summary, initial_share
Lowest in 1988, at 2.98, after a flat decade near 3; it then rose almost without interruption, to 5.11 in 2023. Before 1980
the ratio is dominated by the capital stock the database assumes for 1953 (85 per cent of the 1960 stock), so its fall in
those years says more about that guess than about Korea.
\end{solution}

\begin{solution}{exo:b2:ramsey-model:4}
Poor regions with the same technology and institutions as rich ones that close only a small fraction of the gap each
year, well below the model's 6 per cent; or that do not catch up at all.
\end{solution}

\begin{solution}{exo:b2:ramsey-model:5}
% code: exercises/solutions.py golden
$f'(k) = n + \delta = 0.06$ gives $k_{\mathrm{gold}} = ((1/3)/0.06)^{1.5} = 13.09$, against $k^* = 8.505$. Since $\rho = 0.03 > n = 0.01$, the household discounts the future more than population growth
dilutes it, and does not give up today's consumption to build capital up to the golden rule.
\end{solution}

\begin{solution}{exo:b2:ramsey-model:6}
% code: exercises/solutions.py eigen
$a = \rho - n = 0.02$, $b = c^*f''(k^*)/\theta$ with $f''(k^*) = -\tfrac29 \times 8.505^{-5/3}$; for $\theta = 2$, $b = -0.0048$. The eigenvalues $\tfrac12(0.02 \pm \sqrt{0.0004 + 0.0192})$ are 0.08 and $-0.06$; half-life $\ln 2/0.06 = 11.6$ years. For $\theta = 6$,
$b = -0.0016$, the negative eigenvalue is $\tfrac12(0.02 - \sqrt{0.0004 + 0.0064}) = -0.0312$, and the half-life is 22.2 years.
\end{solution}

\begin{solution}{exo:b2:ramsey-model:7}
% code: exercises/solutions.py steady
Capital falls in the long run, and consumption jumps up at once. $k^* = ((1/3)/0.09)^{1.5} = 7.128$, $c^* = 7.128^{1/3} - 0.06 \times 7.128 = 1.497$, saving rate $0.06 \times 7.128/1.925 = 0.222$. At the moment of the change
capital is still 8.505, now above the new steady state: consumption jumps up onto the new \omterm{def:b2:ramsey-model:saddle-path}{saddle path} and then falls with
capital.
\end{solution}

\begin{solution}{exo:b2:ramsey-model:8}
Today consumption drops, by less than $G$, and capital stays at $k^*$. Until the rise takes effect the old equations govern: with
consumption below $c^*$, capital accumulates, and with $k > k^*$ the return is below $\rho + \delta$, so consumption keeps falling. The path is chosen so that it
reaches the new \omterm{def:b2:ramsey-model:saddle-path}{saddle path} exactly when $G$ rises. From then on capital is run down to $k^*$ and consumption falls to its new
level, $c^* - G$. Households build up capital before the purchases begin and draw on it afterwards.
\end{solution}

\begin{solution}{exo:b2:ramsey-model:9}
% code: exercises/solutions.py euler_growth
$(0.05 - 0.03)/2 = 1$ per cent a year; with $\theta = 1$, 2 per cent.
\end{solution}

\begin{solution}{exo:b2:ramsey-model:10}
With periods of length $\Delta$: $c_t^{-\theta} = e^{-\rho\Delta}e^{r\Delta}c_{t+\Delta}^{-\theta}$, so $(c_{t+\Delta}/c_t)^\theta = e^{(r - \rho)\Delta}$ and $\ln c_{t+\Delta} - \ln c_t = (r - \rho)\Delta/\theta$. Dividing by $\Delta$ and letting $\Delta \to 0$, $\frac{d}{dt}\ln c = \dot c/c = (r - \rho)/\theta$.
\end{solution}

\begin{solution}{exo:b2:ramsey-model:11}
The characteristic polynomial is $\lambda(\lambda - a) + b = \lambda^2 - a\lambda + b$, with roots $\tfrac12(a \pm \sqrt{a^2 - 4b})$. With $b < 0$, $\sqrt{a^2 - 4b} > a$, so the negative root is $\tfrac12(a - \sqrt{a^2 + 4|b|})$, larger in absolute value
when $|b|$ is larger. $|b| = c^*|f''(k^*)|/\theta$ falls with a larger $\theta$ and with a less curved production function, that is with slower diminishing returns.
\end{solution}

\begin{solution}{exo:b2:ramsey-model:12}
The household earns $(1 - \tau)r$ on its saving, so its \omterm{def:b2:intertemporal-choice:euler-equation}{Euler equation} is $\dot c/c = ((1 - \tau)(f'(k) - \delta) - \rho)/\theta$; the lump-sum rebate does not enter it. In the steady state $\dot c = 0$, so $(1 - \tau)(f'(k^*) - \delta) = \rho$, that
is $f'(k^*) = \delta + \rho/(1 - \tau)$, which rises with $\tau$; since $f'' < 0$, $k^*$ falls.
\end{solution}

\begin{solution}{pb:b2:ramsey-model:1}
% code: exercises/solutions.py kestria, initial_saving
\textbf{Part I.}
1. $\dot c/c = (r - 0.04)/2$.
2. $(0.06 - 0.04)/2 = 1$ per cent a year.
3. $1/\theta = 0.5$.
4. Without it the household could choose any starting consumption consistent with the \omterm{def:b2:intertemporal-choice:euler-equation}{Euler equation}, including paths on
which it hoards wealth it never uses or borrows without repaying.
5. $\dot k = k^{0.3} - c - 0.06k$.
\textbf{Part II.}
6. $r = \rho = 0.04$, so $f'(k^*) = 0.10$.
7. $k^* = (0.3/0.1)^{1/0.7} = 3^{1.43} = 4.804$.
8. $y^* = 4.804^{0.3} = 1.601$; $c^* = 1.601 - 0.06 \times 4.804 = 1.313$; wage $0.7 \times 1.601 = 1.121$ thousand kests.
9. $0.06 \times 4.804/1.601 = 0.18$, which is $\alpha\delta/(\rho + \delta) = 0.3 \times 0.06/0.10$.
10. $f'(k) = 0.06$ at $k = 5^{1/0.7} = 9.97$, about twice $k^*$.
11. Not in the model's sense: the gap reflects the household's impatience, which the model takes as part of its welfare;
the market path is the planner's.
\textbf{Part III.}
12. $\begin{pmatrix} 0.04 & -1 \\ -0.0096 & 0 \end{pmatrix}$, with $b = c^*f''(k^*)/2 = 1.313 \times (-0.21 \times 4.804^{-1.7})/2 = -0.0096$.
13. $\tfrac12(0.04 \pm \sqrt{0.0016 + 0.0383}) = 0.1198$ and $-0.0798$.
14. $\ln 2/0.0798 = 8.7$ years.
15. $|b|$ halves; the negative eigenvalue becomes $-0.052$ and the half-life 13.3 years.
16. Households with a larger $\theta$ dislike uneven consumption more, so they accept less of a fall in today's consumption to
build capital fast.
17. Above: shooting gives an initial saving rate of 0.303, since $\theta = 2$ is below $1/0.18 = 5.6$, the case in which the rate falls along the path.
18. 8.7 years means closing 8 per cent of the gap a year, much faster than the slow convergence the evidence on American states
implies.
\textbf{Part IV.}
19. Nothing: the \omterm{def:b2:ramsey-model:modified-golden-rule}{modified golden rule} does not involve $G$.
20. It falls at once by 0.1, from 1.313 to 1.213, and stays there.
21. It does not change: $r = f'(k^*) - \delta = 0.04$ throughout.
22. By less: households spread the cost over the years after the episode by running down capital temporarily.
23. Yes: a tax on capital income lowers the after-tax return, so $k^*$ falls (Exercise 12).
24. Unemployed resources and prices that adjust slowly, which the model rules out by assuming markets always clear.
25. Kestria's steady-state capital is 4.80 per person, its saving rate 18 per cent, and a gap in capital halves in 8.7 years (13.3 with $\theta = 4$).
\end{solution}
